% \sffamily
{\huge \textbf{Abstract}}


The Sun shows a wide range of temporal variations, from a few seconds to decades and even centuries, broadly classified into two classes short\hyp term and \lt. The solar dynamo mechanism is believed to be responsible for these global changes happening in the Sun. Hence, many dynamo models have been proposed to explain the observed behaviour of the Sun. This thesis is primarily focused on studying the \lt\ variation of the Sun and provides  various inputs to the solar dynamo models.

With a renewed interest on the subject several automatic techniques have been developed for extensive data analysis as applied to the long term datasets and presented in this thesis.  This approach provides better consistency and  eliminate human subjectivity, which has been the normal practice in the past. These newly developed techniques are used to extract the umbra and penumbra from sunspots to study area ratio ($q$) and track the sunspot for the measurement of solar differential rotation in the white-light data from the Kodaikanal Solar Observatory (KoSO; 1923\ndash2011), Michelson Doppler Imager (MDI; 1996\ndash2011) and Helioseismic and Magnetic Imager (HMI; 2010\ndash2018). In addition, a completely automatic method to detect the bipolar magnetic regions (BMRs) from the line of sight magnetogram observed using MDI and HMI is also developed. Furthermore, the sunspot area series has been updated for the observation period (1906\ndash2017). We also discovered an issue with the incorrect time stamp in Ca-K data as recorded in the Kodaikanal archive, which we corrected  and oriented the images appropriately.

The ratio of penumbra to umbra area show an increase from 5.5 to 6 as the size of sunspot increases from 100~$\mu$Hem to 2000~$\mu$Hem, it does not show any \lt\ systematic trend as found in Royal Observatory, Greenwich (RGO) photographic results. The solar differential rotation, a  vital parameter in the solar dynamo models, measured in this study is, $\Omega (\theta)= (14.381\pm0.004)-(2.72\pm0.04)\cos^2\theta$, where $\theta$ is the co-latitude. While this study does not show any significant variation on $\Omega$ with time and between the maximum and minimum of activity, it reveals that the bigger sunspots (area$>400~\mu$Hem) give a relatively slower rotation rate than the smaller ones (area$<200~\mu$Hem). The measurement of solar differential rotation using sunspots compared to the spectroscopic observations indicates the existence of an intriguing layer known as the near-surface shear layer (NSSL), which was further verified by a more accurate measurement by helioseismology. The theoretical explanation for NSSL is presented based on the thermal wind balance equation. The idea is based on the fact that, in the top layer of the solar convection zone the temperature profile of the Sun is  independent of latitude; and in addition to that, the sharp fall in temperature causes the thermal wind term to grow. To compensate for the growth of the thermal wind term--the centrifugal term has to increase, leading to the observed NSSL. Based on the observed value of $\Omega (r,\theta)$, $r_c=0.96$~\Rsun\ shows the good agreement with observed NSSL, and it is also hand in hand with the observed estimate of pole equator temperature difference (PETD). Another crucial input for the dynamo model is the observed systematic tilt of the BMRs, particularly in \bl\ types. Based on the thin flux tube models, the tilt induced in BMR depends on the flux tube's rise time, which is a function of the magnetic field it carries. The magnetic field dependence can provide the non-linearity required in kinematic dynamo models to suppress the magnetic field's growth. The distribution of maximum magnetic field (\Bmax) in BMRs is found to be bimodal in both the data sets, peaking at 600~G and 2.2~kG, which turns out to be for BMRs of two classes. The results show an increase in the amplitude of Joy's law for \Bmax$<2$~kG; however, for \Bmax$>2$~kG, it decreases, indicating tilt quenching routinely used in \bl\ models.

The variation of penumbra to umbra area ratio, $q$, observed here will provide constraints in sunspot simulations. In addition, the absence of any difference in the behaviour of small and big spots does not support the idea of the global and local dynamo. Two classes of BMRs observed in the magnetograms further verify this behaviour. The importance of the NSSL is not studied so well in the context of solar dynamo models, but it will be worth waiting to see its significance for understanding solar dynamo. Finally, the indication of tilt quenching presented here needs to be further verified using the more comprehensive data set, including stronger cycles.


\pagebreak
\newpage